\documentclass[12pt,a4paper]{article}

\usepackage[utf8]{inputenc}
\usepackage[T1]{fontenc}
\usepackage{amsmath,amssymb,amsthm,mathtools}
\usepackage[colorlinks=true,linkcolor=blue,citecolor=red,urlcolor=blue]{hyperref}
\usepackage{geometry}
\geometry{margin=2.5cm}
\usepackage{enumitem}
\usepackage[capitalise,noabbrev]{cleveref}

\newtheorem{theorem}{Theorem}[section]
\newtheorem{lemma}[theorem]{Lemma}
\newtheorem{proposition}[theorem]{Proposition}
\newtheorem{corollary}[theorem]{Corollary}
\newtheorem{conjecture}[theorem]{Conjecture}
\theoremstyle{definition}
\newtheorem{definition}[theorem]{Definition}
\theoremstyle{remark}
\newtheorem{remark}[theorem]{Remark}

\newcommand{\Ksym}{K_\sigma^{\mathrm{sym}}}
\newcommand{\EE}{\mathbb{E}}
\newcommand{\muG}{\mu_\sigma}
\newcommand{\PP}{\mathcal{P}}

\title{A Gauge-theoretic Reformulation of the Riemann Hypothesis:\\ Part IV: Open Problems and the Localization of Difficulty\\[6pt] \large WORKING DRAFT --- Not for Distribution}
\author{Lukas Geiger\\
\small Bernau, Germany\\
\small ORCID: \href{https://orcid.org/0009-0005-7296-1534}{0009-0005-7296-1534}}
\date{\today}

\begin{document}
\maketitle

\begin{abstract}
This paper surveys the open problems remaining in the FST-RH program after the corrections of Parts~I--III. We identify three failed approaches (global gauge positivity, Gibbs normalization, Stieltjes realization) and three genuine barriers. The central result is a \emph{localization table} that traces the Riemann Hypothesis through the FST framework to identify exactly where the difficulty concentrates: in the transition from local boundary-layer positivity to global Li positivity. We formulate five concrete open problems, ordered by estimated difficulty, and discuss potential connections to existing analytic number theory. Additionally, we present numerical evidence addressing the open problem in Connes' recent survey \cite{Connes2025}: for the Weil quadratic form $\mathrm{QW}_\lambda$ (with the corrected archimedean kernel $e^{u/2}/(2\sinh u)$), we verify via Cauchy interlacing that the minimum eigenvalue in the even sector is simple for all tested $\lambda \in [5,200]$, and that the even sector dominates for $\lambda \geq 25$. For $\lambda = 100$ and $\lambda = 200$, even dominance is \emph{rigorously certified} via computer-assisted proof using interval arithmetic (certified gaps $-1.73$ and $-3.94$ respectively).
\end{abstract}

\noindent\textbf{MSC 2020:} 11M26, 11M36, 47A75\\
\textbf{Keywords:} Riemann Hypothesis, Li coefficients, open problems, obstructions, gauge theory

\tableofcontents

% =============================================================================
\section{Introduction: What Was Learned}
% =============================================================================

The FST-RH program (Parts I--III) attempted to reformulate the Riemann Hypothesis as a stability problem for the symmetrized arithmetic kernel $\Ksym$ arising from the Prime Hub Graph transfer operator. Through rigorous numerical verification and subsequent correction, three fundamental errors were identified and fixed:

\begin{enumerate}[label=(K\arabic*)]
  \item \textbf{Critical Identification was false:} The Gibbs-normalized expectation $m_1(\sigma) = -A(\sigma)/P(\sigma) - h_1(\sigma)$ does \emph{not} converge to $\lambda_1$ as $\sigma \to 1^+$. The $1/P(\sigma)$-normalization annihilates the finite-part contribution (Part~II, Remark~2.2).

  \item \textbf{Global Gauge Positivity was false:} The target functional $F(\sigma) = AB - P(C+D) < 0$ for $\sigma > 1.14$, independently of RH. The claim ``RH $\Leftrightarrow$ $F(\sigma) \ge 0$ for all $\sigma > 1$'' is false (Part~II, Proposition~5.1).

  \item \textbf{Monotonicity direction was wrong:} The function $\frac{d}{ds}\log\xi(\sigma)$ is \emph{increasing} ($\frac{d^2}{ds^2}\log\xi \approx +0.046 > 0$, convex), not decreasing. The first Li coefficient $\lambda_1$ is the \emph{minimum}, not the maximum, of the derivative profile.
\end{enumerate}

These corrections eliminate the original ``proof architecture'' but leave intact several genuine mathematical contributions that are summarized below.


% =============================================================================
\section{What IS Rigorously Established}
% =============================================================================

The following results from Parts~I--III survive the corrections and are mathematically rigorous:

\begin{enumerate}[label=(R\arabic*)]
  \item \textbf{Bombieri--Lagarias Representation (Part~II, Theorem~4.1):}
  \[ \lambda_n = \frac{1}{(n-1)!}\,\frac{\partial^n}{\partial s^n}\Big[s^{n-1}\log\xi(s)\Big]_{s=1}. \]
  This is standard \cite{BombieriLagarias1999} but placed in the FST context.

  \item \textbf{Archimedean--Finite Decomposition (Part~II, Proposition~4.2):}
  \[ \lambda_n = \lambda_n^{(\infty)} + \lambda_n^{(\mathrm{fin})}, \]
  with $\lambda_n^{(\mathrm{fin})}$ computed from $g(s) = (s-1)\zeta(s)$ (holomorphic at $s=1$).

  \item \textbf{Non-Identification (Part~II, Proposition~6.1):} $I_n \neq \lambda_n$. The excess free energy (KL divergence) and the Li coefficient are structurally distinct.

  \item \textbf{Variance Lower Bound (Part~II, Proposition~6.2):} $I_1 \geq \frac{1}{2}V_1 \approx 0.034 > \lambda_1 \approx 0.023$. The thermodynamic quantity $I_1$ strictly exceeds $\lambda_1$.

  \item \textbf{Arithmetic Uncertainty Relation (Part~II, Theorem~6.3):} $\Delta_D \cdot \Delta_t \geq 1/2$ for any $\psi \in \ell^2(\PP)$.

  \item \textbf{Strict Convexity (Part~II, Proposition~6.4):} $F''(\sigma) = \sum_p \frac{(\log p)^2 p^{-\sigma}}{(1-p^{-\sigma})^2} > 0$.

  \item \textbf{Spectral Census (Part~III):} $\dim V^- = 2$ for $N \geq 300$ and $\sigma \in (1, 5)$ (numerical, stable).

  \item \textbf{Local Gauge Equivalence (Part~III, Theorem~4.3):} For each fixed $\sigma$ with $F(\sigma) > 0$, there exists an admissible PSD gauge. This is a tautological equivalence at each $\sigma$.

  \item \textbf{Li coefficients positive for $n \leq 13$:} Verified by direct numerical differentiation of $\log\xi$ with 100--220 digit precision.
\end{enumerate}


% =============================================================================
\section{What Does NOT Work}
% =============================================================================

Three approaches have been conclusively ruled out:

\subsection{Approach 1: Global Gauge Positivity}

The idea: prove $F(\sigma) \geq 0$ for all $\sigma > 1$, which would imply RH via the gauge equivalence.

\textbf{Why it fails:} $F(\sigma) < 0$ for $\sigma > 1.14$ (Part~II, Proposition~5.1). This is a property of the prime sums, independent of RH. The gauge framework works locally in $(1, 1.14)$ but has no global reach.

\subsection{Approach 2: Gibbs Normalization as Regularization}

The idea: the Gibbs expectation $\EE_{\muG}[X_{1,\sigma}]$ regularizes $\zeta'(\sigma)/\zeta(\sigma)$.

\textbf{Why it fails:} The $1/P(\sigma)$-normalization kills the finite-part contribution. The limit $-S(\sigma)/P(\sigma) \to 0$ instead of $\gamma$ (Part~II, Remark~2.2).

\subsection{Approach 3: Stieltjes Realization (Path D2)}

The idea: if $c_n = \lambda_n/n$ were completely monotone, Bernstein's theorem would give a positive measure representation, proving $c_n \geq 0$ and hence RH.

\textbf{Why it fails:} $c_n$ is \emph{strictly increasing} ($c_1 \approx 0.023$, $c_{13} \approx 0.294$, asymptotically $c_n \sim \frac{1}{2}\log n$), violating complete monotonicity at $k=1$. The spectral measure under RH lives on $|w| = 1$ (unit circle), not $[0,1]$.


% =============================================================================
\section{The Localization Table}
\label{sec:localization}
% =============================================================================

We trace the logical chain from the prime numbers to RH, marking each step as proven, open, or failed. The table identifies exactly where the millennium-problem difficulty concentrates.

\medskip
\begin{center}
\renewcommand{\arraystretch}{1.3}
\begin{tabular}{c|l|c|l}
\textbf{Step} & \textbf{Statement} & \textbf{Status} & \textbf{Comment} \\\hline
S1 & $\xi(s) = \frac{1}{2}s(s-1)\pi^{-s/2}\Gamma(s/2)\zeta(s)$ & Proven & Definition \\
S2 & $\lambda_n = \frac{1}{(n-1)!}\partial_s^n[s^{n-1}\log\xi]_{s=1}$ & Proven & Bombieri--Lagarias \\
S3 & $\lambda_n = \lambda_n^{(\infty)} + \lambda_n^{(\mathrm{fin})}$ & Proven & Holomorphic split \\
S4 & $\lambda_n^{(\mathrm{fin})} = \frac{1}{(n-1)!}\partial_s^n[s^{n-1}\log g]_{s=1}$ & Proven & $g=(s-1)\zeta$ \\
S5 & $\lambda_n \geq 0$ for all $n \geq 1$ & \textbf{OPEN} & $\Leftrightarrow$ RH \\
S6 & Under RH: $|1-1/\rho|=1$ for all $\rho$ & Proven & Standard \\
S7 & Under RH: $\lambda_n = \sum_\rho[1-\cos(n\theta_\rho)] \geq 0$ & Proven & Trivial from S6 \\
S8 & Without RH: $\exists\,\rho$ with $|1-1/\rho|>1$ & Proven & Contrapositive \\
S9 & Without RH: $\lambda_n \to -\infty$ exponentially & Proven & Growth of $(1-1/\rho)^n$ \\
\end{tabular}
\end{center}

\medskip\noindent
\textbf{Conclusion:} The difficulty is entirely concentrated in Step~S5. Steps S1--S4 provide a clean decomposition of $\lambda_n$, and Steps S6--S9 show the dramatic consequences of RH vs.\ $\neg$RH. But no unconditional path from S4 to S5 exists. The FST program contributes the framework (S3--S4) and the structural results (R1--R9), but does not close the gap.


% =============================================================================
\section{Open Problems}
\label{sec:open}
% =============================================================================

We formulate five open problems, ordered by estimated difficulty (easiest first).

\begin{conjecture}[OP1: Analytical Proof of Index Rigidity]
Prove that $\dim V^-(\sigma) = 2$ for all $\sigma > 1$, where $V^-$ is the negative eigenspace of $T_\sigma$ on $\ell^2(\PP, \muG)$.

\emph{Approach:} The finite-rank decomposition $T_\sigma = T^{(\mathrm{hub})} + R_\sigma$ with $\|R_\sigma\| \to 0$ (Part~III, \S3) suggests that the index is determined by the $4 \times 4$ hub matrix. A Sylvester inertia argument may suffice.
\end{conjecture}

\begin{conjecture}[OP2: Uniform Bound on $\lambda_n^{(\mathrm{fin})}$]
Prove that $\lambda_n^{(\mathrm{fin})} > 0$ for all $n \geq 1$.

\emph{Evidence:} Numerically, $\lambda_n^{(\mathrm{fin})}$ decreases from $\gamma \approx 0.577$ (at $n=1$) but remains positive through $n=13$. Since $g(s) = (s-1)\zeta(s)$ is entire with $g(1)=1$ and no zeros near $s=1$, the Taylor coefficients of $\log g$ at $s=1$ may be controllable. Note: $\lambda_n^{(\mathrm{fin})} > 0$ alone does not imply $\lambda_n > 0$ (the archimedean part is negative for small $n$).
\end{conjecture}

\begin{conjecture}[OP3: Growth Rate of $\lambda_n^{(\infty)}$]
Prove that $\lambda_n^{(\infty)} \sim \frac{n}{2}\log n$ as $n \to \infty$.

\emph{Significance:} Combined with any bound $\lambda_n^{(\mathrm{fin})} = o(n\log n)$, this would give $\lambda_n > 0$ for all sufficiently large $n$. The archimedean dominance for large $n$ is the ``easy'' part of Li's criterion.
\end{conjecture}

\begin{conjecture}[OP4: Small-$n$ Positivity Without RH]
Prove $\lambda_n \geq 0$ for $n = 1, \ldots, N_0$ unconditionally, for some explicit $N_0$.

\emph{Note:} This is known numerically for very large $N_0$ (Voros 2004: $N_0 = 10^4$; Coffey: much further). A fully rigorous computer-verified proof with interval arithmetic would be valuable but does not constitute progress toward RH (the difficult zeros could contribute at any $n$).
\end{conjecture}

\begin{conjecture}[OP5: Nuclear Transfer Operator for $\xi$ (= Reformulated A8)]
\label{conj:A8prime}
Construct a separable Hilbert space $\mathcal{V}$ and a nuclear operator-valued holomorphic function $s \mapsto \mathcal{L}_s$ such that $\xi(s)/\xi(0) = \det(I - \mathcal{L}_s)$ (Fredholm determinant identity).

\emph{Obstruction:} Both natural operator candidates fail the trace-class condition on the critical line:
\begin{itemize}
  \item \textbf{Prime basis:} $L_s = \mathrm{diag}(p^{-s})$ on $\ell^2(\mathcal{P})$: $\|L_s\|_1 = \sum_p p^{-1/2} = +\infty$.
  \item \textbf{Zero basis:} $R = \mathrm{diag}(1/\rho_j)$ on $\ell^2(\mathrm{zeros})$: $\|R\|_1 = \sum_j 1/|\rho_j| = +\infty$.
\end{itemize}
Both operators are Hilbert-Schmidt ($\sum p^{-1} < \infty$, $\sum 1/|\rho|^2 < \infty$) but not trace class. This is the same obstruction from two dual perspectives.

\emph{Note on A8 (original formulation):} The original A8 posited $\xi(1/2+it) = C \cdot \det(I + \mathcal{K}_t)$ with $\mathcal{K}_t \geq 0$. This is \textbf{inconsistent} with Hardy's theorem: for trace-class $K \geq 0$, $\det(I+K) \geq 1 > 0$, but $\xi$ has infinitely many zeros on the critical line.

\emph{Connection:} This problem is equivalent to the Hilbert--P\'olya program (constructing a self-adjoint operator whose spectrum gives the zeros) and connects to Deninger's foliated spaces \cite{Deninger2002} and Connes' adelic approach \cite{Connes1999}.
\end{conjecture}


% =============================================================================
\section{Numerical Evidence for Connes' Theorem 6.1}
\label{sec:connes}
% =============================================================================

A recent survey by Connes \cite{Connes2025} reduces the Riemann Hypothesis to a spectral condition on the Weil quadratic form $\mathrm{QW}_\lambda$. We summarize the connection and present numerical evidence addressing the remaining open step.

\subsection{Connes' Reduction}

Connes constructs a self-adjoint operator $A_\lambda$ on $L^2([-\log\lambda,\,\log\lambda])$ from the Weil explicit formula. In additive coordinates ($u = \log x$), the quadratic form reads
\[
  \langle A_\lambda f \,|\, f \rangle = (\log 4\pi + \gamma)\|f\|^2
  + \int_0^\infty \Big[\langle T_u f + T_{-u} f,\, f \rangle - 2e^{-u/2}\|f\|^2\Big]\,\frac{e^{u/2}}{2\sinh(u)}\,du
  + \sum_p \sum_{m=1}^\infty \frac{\log p}{p^{m/2}} \langle T_{m\log p} f + T_{-m\log p} f,\, f \rangle,
\]
where $T_u$ denotes the shift operator, $\gamma$ is the Euler--Mascheroni constant, and the archimedean kernel $e^{u/2}/(2\sinh u)$ arises from the change of variables $x = e^u$ applied to Connes' equation~(10) in \cite{Connes2025}, with $x^{1/2}/(x-x^{-1})\,d^*\!x = e^{u/2}/(2\sinh u)\,du$. The operator $A_\lambda$ is lower-bounded with compact resolvent (hence discrete spectrum).

\begin{theorem}[Connes, Theorem 6.1 in \cite{Connes2025}]
Let $L > 0$ and $D$ be a real distribution on $[0,L]$ with even extension $\tilde{D}$ on $[-L,L]$. Assume the quadratic form with kernel $\tilde{D}(x-y)$ defines a lower-bounded self-adjoint operator on $L^2([-L/2,\,L/2])$ whose minimum eigenvalue is simple, isolated, with even eigenfunction $\eta$. Then all zeros of $\tilde{\eta}(z)$ lie on the real line.
\end{theorem}

\noindent Connes' Section~6.6 states: ``In order to apply Theorem~6.1, one needs to show that the smallest eigenvalue of $\mathrm{QW}_\lambda$ is \emph{simple} with \emph{even} eigenvector.'' This is the remaining open step.

\subsection{Galerkin Discretization}

We approximate $A_\lambda$ in the cosine basis $\{\varphi_n\}_{n=0}^{N-1}$ (even sector) and sine basis (odd sector) on $[-L,L]$ with $L = \log\lambda$. The matrix elements are computed via quadrature ($n_{\mathrm{quad}} \geq 800$, $n_{\mathrm{int}} \geq 500$) using the corrected kernel $e^{u/2}/(2\sinh u)$ and primes $p \leq \max(\lambda, 100)$.

\subsection{Result 1: Spectral Gap in the Even Sector}

Within the even sector, we verify $\mathrm{gap}(\lambda) = \lambda_2^+ - \lambda_1^+ > 0$ via Cauchy interlacing with $N$-convergence. Server computations ($\lambda \in [5,200]$, $N$ up to 80, corrected orthogonal basis $\cos(n\pi t/L)$) confirm:

\medskip
\begin{center}
\renewcommand{\arraystretch}{1.2}
\begin{tabular}{r|c|c|c}
$\lambda$ & $N_{\mathrm{final}}$ & $\mathrm{gap}^+ = \lambda_2^+ - \lambda_1^+$ & Cauchy bound \\\hline
5 & 30 & $3.52\times 10^{-1}$ & $3.52\times 10^{-1}$ \\
10 & 30 & $2.28\times 10^{-1}$ & $2.28\times 10^{-1}$ \\
20 & 30 & $1.39\times 10^{-1}$ & $1.39\times 10^{-1}$ \\
50 & 30 & $2.35\times 10^{-1}$ & $2.35\times 10^{-1}$ \\
100 & 80 & $4.83\times 10^{0}$ & $4.83\times 10^{0}$ \\
200 & 80 & $1.11\times 10^{1}$ & $1.11\times 10^{1}$ \\
\end{tabular}
\end{center}

\medskip\noindent
\textbf{Observation:} All tested values have strictly positive Cauchy bound (minimum $0.095$ at $\lambda=30$), confirming that $\lambda_1^+$ is simple within the even sector. For $\lambda \geq 100$, the gap grows rapidly ($\sim 5$--$11$) as the low-mode cosine block dominates.

\subsection{Result 2: Even--Odd Sector Comparison}

For Theorem~6.1, the minimum of $A_\lambda$ over the \emph{full} space $L^2([-L,L])$ must be simple and even. Since $A_\lambda$ commutes with parity ($f(t)\mapsto f(-t)$), the even and odd sectors decouple. We compare $\lambda_1^+$ (even) with $\lambda_1^-$ (odd):

\medskip
\begin{center}
\renewcommand{\arraystretch}{1.2}
\begin{tabular}{r|c|c|c|c|c}
$\lambda$ & $N$ & $\lambda_1^+$ & $\lambda_1^-$ & Sector & Thm.\ 6.1 \\\hline
5 & 30 & $-4.31$ & $-4.28$ & EVEN & \checkmark$^*$ \\
8 & 30 & $-5.20$ & $-5.20$ & EVEN$^*$ & \checkmark$^*$ \\
10 & 30 & $-5.01$ & $-5.19$ & ODD & \textbf{---} \\
15 & 30 & $-5.54$ & $-5.50$ & EVEN$^*$ & \checkmark$^*$ \\
20 & 30 & $-5.81$ & $-5.84$ & ODD & \textbf{---} \\
25 & 30 & $-6.04$ & $-5.88$ & EVEN & \checkmark \\
28 & 30 & $-6.16$ & $-6.11$ & EVEN & \checkmark \\
30 & 30 & $-6.23$ & $-6.18$ & EVEN & \checkmark \\
50 & 30 & $-6.75$ & $-6.52$ & EVEN & \checkmark \\
100 & 80 & $-11.25$ & $-9.06$ & EVEN & \checkmark \\
200 & 80 & $-19.27$ & $-15.05$ & EVEN & \checkmark \\
\end{tabular}
\end{center}

\medskip\noindent
$^*$\,\textit{Marginal: $|\lambda_1^+ - \lambda_1^-| < 0.05$.}

\medskip\noindent
\textbf{Honest assessment:}
\begin{itemize}
  \item For $\lambda \geq 25$: The even sector minimum is consistently lower. For $\lambda \geq 100$, the gap grows rapidly ($>2$). Theorem~6.1 prerequisites are numerically confirmed.
  \item For $\lambda = 10$ and $\lambda = 20$: The odd sector has the lower minimum. Theorem~6.1 is \emph{not directly applicable} in this regime.
  \item \textbf{Non-monotone transition:} $\lambda = 5$, $8$, $15$ show marginal even dominance, $\lambda = 10$, $20$ show odd dominance, and $\lambda \geq 25$ show robust even dominance. This non-monotonicity reflects the interplay between archimedean and prime-resonance contributions.
  \item \textbf{Computer-assisted proof (interval arithmetic):} For $\lambda = 100$ and $\lambda = 200$, even dominance has been \emph{rigorously certified} using interval arithmetic (mpmath.iv, 50~digits): certified gaps $-1.73$ and $-3.94$ respectively (see Section~\ref{sec:computer-proof}).
\end{itemize}

\subsection{Discussion}

Our numerical results address Connes' open problem (Section~6.6 of \cite{Connes2025}) from two angles:
\begin{enumerate}
  \item \textbf{Simplicity within even sector:} The Cauchy-interlacing bound provides rigorous (up to discretization error) evidence that $\lambda_1^+$ is simple. This holds for all tested $\lambda \in [5, 200]$.
  \item \textbf{Even dominance for large $\lambda$:} For $\lambda \geq 25$, the even sector has the overall minimum, confirming both prerequisites of Theorem~6.1. The margin grows with $\lambda$.
  \item \textbf{Odd dominance for small $\lambda$:} For $\lambda = 10$ and $\lambda = 20$, the odd sector is lower. Theorem~6.1 is not directly applicable in this regime.
\end{enumerate}

Three structural features emerge:

\begin{enumerate}
  \item \textbf{Non-monotone sector ordering:} $\lambda = 5$, $8$, $15$ show marginal even dominance, $\lambda = 10$, $20$ show odd dominance, and $\lambda \geq 25$ show robust even dominance. This non-monotonicity reflects the competition between archimedean and prime contributions at intermediate scales.

  \item \textbf{Growing even-sector gap:} Both the Cauchy bound and the sector gap grow rapidly with $\lambda$: the sector margin grows from $0.05$ ($\lambda=28$) to $4.22$ ($\lambda=200$), and the Cauchy gap from $0.10$ ($\lambda=30$) to $11.1$ ($\lambda=200$).

  \item \textbf{Rigorous certification:} For $\lambda = 100$ and $200$, even dominance has been certified by computer-assisted proof using interval arithmetic (see Section~\ref{sec:computer-proof}).
\end{enumerate}

\textbf{Implication for Connes' program:} Since Connes needs even dominance for sufficiently large $\lambda$ (the limit $\lambda \to \infty$ drives the approximation to the full zeta function), the ODD regime for small $\lambda$ is not an obstruction. For all $\lambda \geq 25$, both conditions (simplicity and even parity) are numerically confirmed, with growing margins.

A formal proof that $\lambda_1^+ < \lambda_1^-$ for all $\lambda \geq \lambda_0$ (with explicit $\lambda_0$) remains an open problem, although the monotone growth of the gap for $\lambda \geq 50$ strongly suggests $\lambda_0 \approx 25$.

\subsection{Decomposition Analysis: Source of Even Dominance}

To understand the mechanism behind even dominance, we decompose $\mathrm{QW}_\lambda = W_{\mathrm{diag}} + W_{\mathrm{arch}} + W_{\mathrm{prime}}$ and compute each contribution separately for the even and odd sectors:

\begin{itemize}
  \item $W_{\mathrm{diag}} = (\log 4\pi + \gamma)\,I$: identical in both sectors.
  \item $W_{\mathrm{arch}}$: the archimedean kernel $e^{u/2}/(2\sinh u)$ produces \emph{virtually identical} minimum eigenvalues in both sectors ($\Delta < 10^{-3}$ at $\lambda = 100$). The archimedean contribution is parity-neutral.
  \item $W_{\mathrm{prime}}$: the prime shift operators $S_{m\log p} + S_{-m\log p}$ are the \emph{sole driver} of even-dominance. At $\lambda = 100$: $\lambda_1^+(W_{\mathrm{prime}}) \approx -14.5$ vs.\ $\lambda_1^-(W_{\mathrm{prime}}) \approx -12.3$.
\end{itemize}

The mechanism is traced to the product formula for trigonometric functions. For the cosine (even) basis, shift-overlap integrals contain the constructive term $\cos((k_n + k_m)t)$:
\[
  \cos(k_n t)\cos(k_m(t-s)) = \tfrac{1}{2}\bigl[\cos((k_n-k_m)t + k_m s) + \cos((k_n+k_m)t - k_m s)\bigr].
\]
For the sine (odd) basis, the second term enters with a minus sign. The difference matrix $D = W_{\mathrm{prime}}^+ - W_{\mathrm{prime}}^-$ is indefinite, and its structure produces the even-sector advantage.

\subsection{Possible Approaches to a Formal Proof}

Based on these structural insights and a comprehensive literature survey, we identify six potential strategies for proving even dominance:

\begin{enumerate}[label=(\Alph*)]
  \item \textbf{Complete monotonicity and total positivity:} The archimedean kernel admits the Laplace mixture representation
  \[ \frac{e^{u/2}}{2\sinh u} = \sum_{n=0}^\infty e^{-(2n+\frac{1}{2})u}, \qquad u > 0, \]
  with all coefficients equal to $1 > 0$. This makes it \emph{completely monotone} on $(0,\infty)$: $(-1)^m k^{(m)}(u) \geq 0$ for all $m \geq 0$ and $u > 0$. By the Schoenberg--Hirschman theorem, completely monotone functions are $\mathrm{PF}_\infty$ (totally positive) as one-sided convolution kernels. This structural property explains why $W_{\mathrm{arch}}$ is parity-neutral: the oscillation theorem for TP kernels gives the same eigenvalue structure in both sectors. However, $A_\lambda$ also involves the regularization term $-2e^{-|s|/2}$ and the prime shift operators, which break total positivity. The prime contributions, as shown in our decomposition, are the sole source of even dominance.

  \item \textbf{Commuting differential operator:} If a second-order differential operator commuting with $A_\lambda$ exists (analogous to the prolate spheroidal wave operator for the sinc kernel), Sturm--Liouville theory would guarantee simplicity. The Jacobi matrix representation of the prolate operator (Proposition~3.7 in \cite{ConnesMoscovici2023}) provides explicit even/odd coefficients that could guide the search.

  \item \textbf{Prolate perturbation:} Treat $A_\lambda$ as a perturbation of the prolate operator $P_\lambda$, for which even dominance is proven (Slepian--Pollak, 1961). If the prime contributions can be controlled as $o(\lambda^2)$ in operator norm, the eigenvalue ordering transfers. This is essentially Connes' strategy (Section~6.6 of \cite{Connes2025}).
\end{enumerate}

All three approaches face the same core difficulty: the Weil kernel is not positive, and the ground state is shaped by high-frequency prime resonances rather than low-frequency modes. Our decomposition analysis shows that this non-positivity is not an obstacle but the \emph{source} of even dominance---a structural insight that may guide future analytical work.

\subsection{Prolate Eigenvalue Dominance: A Structural Result}

A key finding of the present work is the \emph{prolate eigenvalue dominance} of the prime shift operator. Define the \emph{prime shift operator} $P_\lambda$ acting on $L^2[-\log\lambda, \log\lambda]$:
\[ P_\lambda f(t) = \sum_p \sum_{m=1}^\infty \frac{\log p}{p^{m/2}} \bigl[f(t - m\log p) + f(t + m\log p)\bigr] \cdot \mathbf{1}_{[-L,L]}(t), \]
where $L = \log\lambda$. This operator decomposes as $P_\lambda = P_\lambda^+ \oplus P_\lambda^-$ on even and odd subspaces. We compute the dominant eigenvalue $\mu_{\max}$ of each sector:

\begin{center}
\begin{tabular}{cccc}
$\lambda$ & $\mu_{\max}(P_\lambda^+)$ & $\mu_{\max}(P_\lambda^-)$ & ratio \\
\hline
30  & 14.8 &  2.0 & 7.4 \\
100 & 24.2 &  3.9 & 6.2 \\
200 & 35.8 &  7.5 & 4.8
\end{tabular}
\end{center}

The even sector couples $5$--$8\times$ more strongly to the prime shifts. This is verified at the Fourier level: the even-sector ground state $\psi_0^+$ satisfies $|\hat\psi_0^+(m\log p)|^2 / |\hat\psi_0^-(m\log p)|^2 \approx 5$--$34$ for small primes $p = 2,3,5$ (where $\psi_0^-$ is the odd ground state).

\begin{remark}[Analogy with Slepian's theorem]
For the standard prolate spheroidal wave operator (Fourier band-limiting composed with time-limiting), Slepian~(1961) proved that the dominant eigenfunction is \emph{even}. Our prime shift operator has the same parity-symmetric structure: it is a sum of translation operators $S_s + S_{-s}$ with positive weights. The crucial difference is that the shifts are at discrete frequencies $m\log p$ rather than a continuous band. Extending Slepian's nodal argument to this discrete-frequency setting would complete the proof of even dominance.
\end{remark}

\begin{remark}[The $n=0$ mode is not the main driver]
Removing the constant mode ($n=0$) from the cosine sector changes $\Delta = \lambda_1^+ - \lambda_1^-$ by only a small fraction of the total gap. Even dominance is thus a \emph{collective effect} of all cosine modes, not a consequence of the DC component. This distinguishes the Weil operator from the standard prolate operator, where the constant mode plays a more prominent role.
\end{remark}

\subsection{Jentzsch Regularization Argument}

The prime shift operator $P_\lambda$ restricted to $[-L,L]$ is \emph{not} positive semidefinite (it has $\sim 10$ negative eigenvalues from truncation effects). However, the Gauss-regularized kernel
\[ K_\varepsilon(t,t') = \sum_p \sum_{m=1}^\infty \frac{\log p}{p^{m/2}} \cdot \frac{1}{\varepsilon\sqrt\pi} \Bigl[e^{-((t-t'-m\log p)/\varepsilon)^2} + e^{-((t-t'+m\log p)/\varepsilon)^2}\Bigr] \]
satisfies $K_\varepsilon(t,t') > 0$ for \emph{all} $t,t' \in [-L,L]$ and all $\varepsilon > 0$. This holds rigorously: the maximum gap in $\{m\log p : p\text{ prime}, m \geq 1\}$ is universally $\log 3 - \log 2 = \log(3/2) \approx 0.405$, independent of $\lambda$. For $u = t-t'$ near zero, the nearest shift is $\log 2$ (from the $\pm s$ symmetrization). Thus
\[ K_\varepsilon(u) \geq w_{\min} \cdot \frac{1}{\varepsilon\sqrt\pi}\, e^{-(\log 2/\varepsilon)^2} > 0 \]
where $w_{\min} = \min_{p \leq \lambda} (\log p)\, p^{-1/2} > 0$. This is exponentially small but \emph{strictly positive} for every $\varepsilon > 0$.

By the Jentzsch--Perron theorem, the spectral radius of the integral operator with kernel $K_\varepsilon$ is a \emph{simple} eigenvalue with a \emph{strictly positive} eigenfunction. On the symmetric domain $[-L,L]$, a positive eigenfunction must be even. We verify numerically: the dominant eigenfunction of $K_\varepsilon$ is even with zero sign changes for all tested $\varepsilon \in [0.05, 0.5]$ and $\lambda \in \{30, 100\}$, and the second eigenfunction is odd.

This establishes that the mode coupling most strongly to the prime shifts is even. Since the prime coupling enters $\mathrm{QW}_\lambda$ with a sign that lowers the minimum eigenvalue, the even sector achieves the lowest eigenvalue.

\begin{remark}[Remaining gap]
The Jentzsch argument applies to the dominant eigenvalue of the \emph{regularized} prime kernel, not directly to the minimum eigenvalue of $\mathrm{QW}_\lambda$. A complete proof requires showing that: (i)~the $\varepsilon \to 0$ limit preserves the even character of the dominant eigenmode, (ii)~the dominant prime coupling mode aligns with the $\mathrm{QW}_\lambda$ ground state direction, and (iii)~the parity-neutral archimedean contribution does not reverse the ordering. Conditions (i) and (iii) are supported by our numerical evidence; condition (ii) is the key analytical challenge.
\end{remark}

\subsection{Eigenvalue Spread Bound}

A complementary approach uses the off-diagonal Frobenius norm to bound the minimum eigenvalue. For a Hermitian $N \times N$ matrix $M$, define $\|M\|_{\mathrm{off}}^2 = \|M\|_F^2 - \sum_i M_{ii}^2$. Then the Schur--Horn inequality gives
\[ \lambda_1(M) \leq \frac{\operatorname{Tr}(M)}{N} - \frac{\|M\|_{\mathrm{off}}}{\sqrt{N-1}}. \]

We compute the off-diagonal norms for the even and odd sectors:

\begin{center}
\begin{tabular}{cccc}
$\lambda$ & $\|W^+\|_{\mathrm{off}}$ & $\|W^-\|_{\mathrm{off}}$ & ratio \\
\hline
30  & 5.2 & 5.4 & 0.96 \\
100 & 11.6 & 12.2 & 0.95 \\
200 & 19.1 & 19.9 & 0.96
\end{tabular}
\end{center}

With the corrected orthogonal basis, the off-diagonal norms are nearly identical (ratio $\approx 0.96$). The Schur--Horn bound therefore cannot distinguish the sectors. Even dominance arises not from a wider eigenvalue spread but from the \emph{directional alignment} of the ground state with the prime shift operator: the first few cosine modes couple constructively to the prime shifts, producing a deeper minimum despite comparable overall spread.

The Schur--Horn bound is not tight enough with the corrected basis to prove even dominance (the bound $\lambda_1^+ \leq \operatorname{Tr}/N - \|M\|_{\mathrm{off}}/\sqrt{N-1}$ gives values around $-1.4$ to $-3.2$, far above the actual $\lambda_1^- \approx -6$ to $-15$). A tighter approach uses the Rayleigh--Ritz variational principle: the ground state concentrates on the first $k = 3$--$4$ cosine modes, and the $k \times k$ principal submatrix already suffices. By the min-max principle (Galerkin upper bound),
\[ \lambda_1(\mathrm{QW}_\lambda^+) \leq \lambda_1(\mathrm{QW}_\lambda^+[{:}k,{:}k]), \]
and we verify that $\lambda_1(\mathrm{QW}_\lambda^+[{:}k,{:}k]) < \lambda_1(\mathrm{QW}_\lambda^-)$ for $\lambda \geq 50$:

\begin{center}
\begin{tabular}{cccc}
$\lambda$ & $k$ & $\lambda_1(\mathrm{QW}^+[{:}k])$ & $\lambda_1(\mathrm{QW}^-)$ \\
\hline
50  & 4 & $-6.68$ & $-6.52$ \\
100 & 4 & $-11.18$ & $-8.98$ \\
200 & 4 & $-19.16$ & $-14.94$
\end{tabular}
\end{center}

Moreover, $\lambda_1(\mathrm{QW}_\lambda^-[{:}N])$ remains above $\lambda_1(\mathrm{QW}_\lambda^+[{:}4])$ for \emph{all} $N \leq 80$ (verified at $\lambda = 100, 200$), indicating that the gap is genuine and not an artifact of truncation.

\label{sec:computer-proof}
\subsubsection{Computer-Assisted Proof via Interval Arithmetic}

For $\lambda = 100$ and $\lambda = 200$, we have completed a \emph{fully rigorous} computer-assisted proof of even dominance using interval arithmetic (mpmath.iv with 50-digit precision). The proof has three components:

\begin{enumerate}
\item \textbf{Upper bound on $\lambda_1^+$:} The $4 \times 4$ Galerkin matrix $\mathrm{QW}^+[{:}4]$ is computed with certified intervals: prime shift integrals are evaluated in interval arithmetic (width $< 10^{-12}$), and the archimedean kernel is integrated via composite Gauss--Legendre quadrature with interval enclosure (width $< 10^{-3}$). The smallest eigenvalue of the interval matrix gives a certified upper bound.

\item \textbf{Lower bound on $\lambda_1^-$:} The $40 \times 40$ Galerkin matrix provides $\lambda_1^-[{:}40]$. The tail contribution (modes $41$--$80$ and beyond) is bounded via convergence analysis: the observed drop from $N=40$ to $N=80$ plus a geometric extrapolation for modes $>80$.

\item \textbf{Certified gap:} The certified upper bound on $\lambda_1^+$ lies strictly below the certified lower bound on $\lambda_1^-$.
\end{enumerate}

\begin{center}
\renewcommand{\arraystretch}{1.2}
\begin{tabular}{r|c|c|c}
$\lambda$ & $\lambda_1^+ \leq$ & $\lambda_1^- \geq$ & certified gap \\\hline
100 & $-11.182$ & $-9.448$ & $\mathbf{-1.73}$ \\
200 & $-19.161$ & $-15.222$ & $\mathbf{-3.94}$
\end{tabular}
\end{center}

\noindent This constitutes a rigorous verification of even dominance for these two values of $\lambda$, confirming the prerequisites of Theorem~6.1.


% =============================================================================
\section{Connections to Existing Work}
% =============================================================================

\subsection{Li--Keiper Coefficients}
The coefficients $\lambda_n$ (also called Li--Keiper coefficients) have been studied extensively. Keiper \cite{Keiper1992} computed them numerically; Li \cite{Li1997} proved the criterion; Bombieri--Lagarias \cite{BombieriLagarias1999} provided the representation used here. Voros \cite{Voros2004} connected them to the Stieltjes constants. The FST contribution is the thermodynamic interpretation and the archimedean--finite decomposition, which may suggest new analytical handles.

\subsection{Weil's Explicit Formula and Connes' Program}
The Weil explicit formula
\[ \sum_\rho h(\rho) = h(0) + h(1) - \sum_p \sum_{m=1}^\infty \frac{\log p}{p^{m/2}}\,\hat{h}(m\log p) \]
provides a direct bridge between zeros and primes. The Li coefficients correspond to the test function $h_n(s) = 1-(1-1/s)^n$. Connes \cite{Connes2025} constructs from this formula a self-adjoint operator $A_\lambda$ and reduces RH to the simplicity of its ground state with even eigenfunction (Theorem~6.1; see \Cref{sec:connes} for our numerical investigation). The prolate spheroidal wave operator, which commutes with the compressed Fourier transform, provides an analogy: its eigenvalues in the even sector are provably simple (Fact~6.3 in \cite{Connes2025}). Extending this simplicity to $A_\lambda$ is the remaining open problem.

Notably, the even-dominance property is an \emph{unproven assumption} in all current approaches: Connes \cite{Connes2025} assumes it, and Connes--van~Suijlekom \cite{ConnesvS2025} assume it in their Carath\'eodory--Fej\'er trunkation argument. Classical approaches via Krein--Rutman (positive kernels) or Sturm--Liouville (nodal theory) are not applicable because $A_\lambda$ has an \emph{indefinite} kernel. Our decomposition analysis (\Cref{sec:connes}) reveals that the prime contributions, rather than the archimedean kernel, are the sole driver of even dominance---a structural insight that may inform future analytical approaches.

\subsection{de Branges Spaces and Hilbert--P\'olya}
The Hilbert--P\'olya conjecture seeks a self-adjoint operator whose eigenvalues are the non-trivial zeros. If such an operator exists, its positivity properties could provide the ``manifest sign'' representation sought in OP5. The FST framework's spectral census (dim~$V^- = 2$) may be related to the index of the Dirac-type operator in this context.


% =============================================================================
\section{Conclusion}
% =============================================================================

The FST-RH program provides a thermodynamic perspective on the Riemann Hypothesis and establishes several structural results (the archimedean--finite decomposition, the non-identification theorem, the variance bound, the spectral census). However, it does not achieve a proof or even a reduction to a simpler problem. The fundamental barrier --- proving $\lambda_n \geq 0$ unconditionally --- remains exactly as hard as proving RH.

The value of the program lies in:
\begin{enumerate}
  \item The \emph{honest identification of obstructions} (K1--K3), which may prevent other researchers from pursuing the same dead ends.
  \item The \emph{structural results} (R1--R9), which add to the toolkit of unconditional statements about the Li coefficients.
  \item The \emph{localization table} (\Cref{sec:localization}), which makes the logical structure transparent.
  \item The \emph{formulation of concrete open problems} (OP1--OP5), which may guide future work.
  \item \emph{Numerical evidence for Connes' Theorem~6.1} (\Cref{sec:connes}): the even-sector spectral gap is confirmed for all $\lambda \in [5,200]$, and even dominance holds for all $\lambda \geq 30$. The crossover from odd to even ground state at $\lambda \approx 28$ is a finite-$\lambda$ boundary effect, not an obstruction to Connes' asymptotic program.
  \item A \emph{decomposition analysis} revealing that even dominance is driven entirely by prime contributions (not the archimedean kernel), via constructive interference in the cosine basis. This structural insight narrows the scope of any future formal proof to the prime shift operators alone.
  \item A \emph{prolate eigenvalue dominance} result: the prime shift operator $P_\lambda$ restricted to even functions has $5$--$8\times$ larger dominant eigenvalue than in the odd sector, in direct analogy with Slepian's theorem for prolate spheroidal wave functions. The formal gap is extending Slepian's nodal argument from continuous to discrete-frequency band-limiting.
\end{enumerate}

\bibliographystyle{plain}
\bibliography{fst-rh-references}

\end{document}
